\section{DeFi Experiment}
\label{sec:defi-experiment}

We now port the DA-BiGRU-CNN of
\Cref{sec:methodology-abstract,sec:methodology-defi} to Aave V3 and
Compound V3 USDC supply markets and evaluate a forecast-driven 4-factor
MCDM allocator against a reactive EMA-smoothed baseline.  The headline
test was pre-registered as $H_1{:}~\Delta\textsc{Sharpe}\ge 0.2$
(\Cref{sec:introduction}); the publication protocol commits to reporting
the binding result whether or not $H_1$ is rejected.  As anticipated
under the H0-publishable framing, the test window does \emph{not} reject
$H_0$; we present both the forecaster diagnostics that drive that
outcome and the design ablation that confirms the allocator itself is
calibrated rather than mis-specified.

\subsection{Data and splits}
\label{sec:defi-data}

The training panel comprises $\DataHours$ hourly bars from $\DataStart$
to $\DataEnd$ at $\DataCoverage$ cross-protocol coverage, joined via the
hourly resampling pipeline of \texttt{data/clean.py}.  Aave V3 rates are
fetched from the official subgraph
\cite{aavecompoundcrosschain2025,gudgeon2020loanable}; Compound V3 base
rates are reconstructed from on-chain view calls because the Messari
subgraph does not index the base Comet market
(see \Cref{sec:limitations} for details on the data-engineering
constraints).  The split is the standard chronological 70/15/15 train /
validation / test partition; the test fold covers the final four months,
which (per the regime structure of \Cref{sec:bg-defi-mechanics})
straddles the 2025-Q4 inversion and the 2026-Q2 partial window in which
Aave reverts to dominating.  This is the most adversarial split available
on the panel and was chosen deliberately to test transferability under
distribution shift rather than under in-sample stationarity.

\subsection{Forecaster training and diagnostics}
\label{sec:defi-forecaster}

The DA-BiGRU-CNN of \Cref{sec:methodology-defi} ($\ModelParams$
parameters, $T=\SeqLen$ hour context, $H=\Horizon$ hour horizon) is
trained with the composite loss of \Cref{eq:composite-loss} under the
same $(\alpha,\beta,\gamma)=(0.4,0.5,0.1)$ schedule used in the LOB
specialisation \cite{solovev2026lob}.  On the validation fold the
network attains weighted-Pearson $0.747$ on Aave and $0.289$ on
Compound, with directional accuracy $0.739$ and a positive
$R^{2}=0.138$ on the Compound branch --- the latter beats a constant
mean predictor, confirming non-trivial signal extraction at the
in-distribution level.

On the held-out test fold the picture inverts.  Weighted-Pearson is
$-0.07$ on Aave and $+0.603$ on Compound, and directional accuracy
drops to $0.351$, below the $0.5$ random-coin baseline.  The 2025-Q3
$\to$ Q4 regime crossover documented in \Cref{sec:bg-defi-mechanics} ---
the Aave-pays-more share rising from $39.9\%$ to $56.3\%$, with the
median spread inverting from $-0.22$\,pp to $+0.18$\,pp --- coincides
with this drop and is the most plausible attribution.  The validation
fold sits inside the pre-inversion regime; the test fold straddles the
inversion and a return to Compound dominance in 2026-Q1.  The
sub-chance directional accuracy is the empirical signature of
forecasting a stationary process whose underlying mean has shifted: the
model commits, with high confidence in well-calibrated weighted-Pearson
terms, to the wrong side of the cross-protocol comparison.

\subsection{Headline H1 test}
\label{sec:defi-h1}

The forecast vector feeds the 4-factor MCDM allocator of
\Cref{sec:bg-mcdm} (APY, Risk, Cost, Stability), with rebalance gated
by a hysteresis threshold $\theta=0.05$ on the criteria-score delta.
\Cref{tab:h1-headline} reports the predictive variant against the
EMA-smoothed reactive baseline over the 4-month test window.

\begin{table}[t]
  \centering
  \caption{Headline H1 test on the 4-month test window. The predictive
    variant under-performs the reactive EMA baseline by 8 bp net APY;
    the bootstrap 95\% CI on the monthly-Sharpe difference lies
    strictly below zero, so $H_1$ is not rejected.}
  \label{tab:h1-headline}
  \begin{tabular}{lcccc}
    \toprule
    Strategy & Net APY & Sharpe & $n_{\text{rebal}}$ & Gas (\$) \\
    \midrule
    PredictiveMCDM & \PredictiveAPY & \PredictiveSharpe & 2 & \PredictiveGas \\
    MCDMEMA        & \EMAAPY        & \EMASharpe        & 2 & \EMAGas \\
    \bottomrule
  \end{tabular}
\end{table}

A 1000-resample paired monthly-Sharpe bootstrap (PredictiveMCDM minus
MCDMEMA on per-month Sharpe pairs) returns
$\Delta\textsc{Sharpe}=\SharpeDelta$ with 95\% confidence interval
$[\BootstrapCILow,\,\BootstrapCIHigh]$ and bootstrap $p$-value
$\BootstrapPvalue$.  Both endpoints of the CI lie below the
$H_1$ threshold of $+0.2$; the verdict is therefore: \emph{fail to
reject $H_0$}.  The wide CI magnitude is a four-month-window artifact
(only four paired observations) and is the same caveat documented in
the pre-registration; with a longer test fold the same point estimate
would carry a tighter interval, but the sign of the contrast would not
change.

\subsection{Hysteresis-sweep ablation}
\label{sec:defi-ablation}

A natural concern with the H1 result is that the default hysteresis
threshold $\theta=0.05$ is mis-calibrated.  Ablation~\#16 sweeps
$\theta\in\{0.01, 0.02, 0.05, 0.10, 0.20\}$ on PredictiveMCDM with all
other parameters fixed (\Cref{tab:ablation-16}).  Increasing rebalance
frequency from $0$ to $7$ \emph{hurts} risk-adjusted return: net APY
falls from $\AblationSixteenDAPY$ at zero rebalances to
$\AblationSixteenAAPY$ at seven, even though Sharpe rises into the
$\AblationSixteenASharpe$--$\AblationSixteenCSharpe$ range.  The default
$\theta=0.05$ sits at the knee of the curve, with one rebalance,
$\AblationSixteenCSharpe$ Sharpe, and $\AblationSixteenCAPY$ APY ---
gas-light and risk-adjusted, exactly the calibration the
pre-registration prescribed.  The ablation does not change the H1
verdict, but it does establish that the design is correctly tuned: the
gap to EMA is not a hysteresis-knob artifact.

\begin{table}[t]
  \centering
  \caption{Ablation~\#16: hysteresis-threshold sweep on PredictiveMCDM.
    Default $\theta=0.05$ (row C) sits at the knee.}
  \label{tab:ablation-16}
  \begin{tabular}{lccc}
    \toprule
    $\theta$ & Net APY & Sharpe & $n_{\text{rebal}}$ \\
    \midrule
    0.01  & \AblationSixteenAAPY & \AblationSixteenASharpe & \AblationSixteenARebalances \\
    0.02  & \AblationSixteenBAPY & \AblationSixteenBSharpe & \AblationSixteenBRebalances \\
    0.05  & \AblationSixteenCAPY & \AblationSixteenCSharpe & \AblationSixteenCRebalances \\
    0.10  & \AblationSixteenDAPY & \AblationSixteenDSharpe & \AblationSixteenDRebalances \\
    0.20  & \AblationSixteenEAPY & \AblationSixteenESharpe & \AblationSixteenERebalances \\
    \bottomrule
  \end{tabular}
\end{table}

In summary, the architecture extracts in-sample structure on both
protocols (validation weighted-Pearson $0.747$ on Aave, directional
accuracy $0.739$); the allocator is correctly calibrated (one rebalance
at the design $\theta$, gas-light); and yet the headline H1 fails on
the binding test fold, traceable to a documented regime shift. The
discussion in \Cref{sec:discussion} unpacks what this means for
cross-domain transfer.
